\section{Contact with the central section and positive offsets}
\label{sec:contact-positive-offsets}

We begin with two regions in which \(\Delta_p\) has at most two zeros: the contact scale
\(q_0=\widehat q_0q_1^2\) with bounded \(\widehat q_0\), and positive
offsets with small \(q_1/\sqrt{q_0}\). In the first, a scattering
comparison controls the derivative ratio on a compact range of central
heights, while the paired identity controls the two ends of the passing
interval. In the second, the Dulac map of an interior saddle determines
the crossing direction of the logarithmic slope ratio near the
separatrices; exterior contraction controls the remaining heights.

In this section and in
\cref{sec:recurring-ranges,sec:primary-regimes,sec:quadratic-regimes,sec:direct-negative-annulus-buffer,sec:direct-negative-root-regimes,sec:sharpness},
the height \(z_*\) of the fixed section, the exterior windows, the common
cone, and the landing level \(m_{\rm land}\) are those chosen in the first
row of \cref{tab:direct-constant-order}. All scaled thresholds and parameter
bounds below are chosen after these fixed data.

The constant \(C_{\rm cone}\) in \eqref{eq:direct-tail-first-slope} is
independent of the contact width and the later coefficient margins and
regional cuts.
The continuation constants and higher exterior derivatives may depend
on these cuts and on \(z_*\); they enter only the final parameter
thresholds, with the sections fixed. The section requirements are justified
in the proof of \cref{prop:direct-uniform-displacement} and the derivation
of \eqref{eq:sharpness-lambda-range}.
